interest except modulo torsion, i.e. after tensoring with @@M0@@. At the same time, this example shows that only after tensoring with @@M1@@ is @@M2@@ covariant with respect to proper morphisms, even when restricted to closed immersions of smooth quasi-projective schemes. Indeed, for @@M3@@ such that @@M4@@ is not an isomorphism, there is in @@M5@@ an integral closed subscheme @@M6@@ such that @@M7@@ does not lie in @@M8@@, where @@M9@@. We may take @@M10@@ of minimal dimension, so that @@M11@@. With @@M12@@ the singular locus of @@M13@@ and @@M14@@, @@M15@@, we conclude that @@M16@@ does not lie in @@M17@@ (since @@M18@@ is surjective and compatible with the augmentations, and its kernel is @@M19@@). Now consider the inclusion
@@M20@@
which is a closed immersion of codimension @@M21@@ of smooth quasi-projective schemes over @@M22@@, \uline{such that}
@@M23@@
\uline{is not compatible with the} @@M24@@-\uline{filtrations} (\uline{up to a shift by} @@M25@@), since @@M26@@ does not lie in @@M27@@.

\sgaStableTarget{sga6.e14.para.4-7}\textbf{4.7.} Keeping the notation of 4.5, observe now that (Exp. IX) shows that the cited theorem of HODGKIN is equivalent to the existence of an isomorphism
@@M28@@
the homomorphism @@M29@@ being a priori surjective, with kernel the ideal generated by the @@M30@@, where @@M31@@ ranges over the invertible Modules on @@M32@@ associated with the representations @@M33@@ (which in fact give all the invertible modules on @@M34@@ [5, Exp. 15]). On the other hand, an analogous argument shows that
@@M35@@
is also surjective and that its kernel is the ideal generated by the elements of @@M36@@ defined by the representations @@M37@@, thus in fact by all of @@M38@@.\SGAAuthorFootnote{*}{Cf. IX 3.4, 2.2.} This shows that @@M39@@ reduces to @@M40@@ if and
